% !Mode:: "TeX:UTF-8"

%? ************************************************************* %?
%?                                                               %?
%?                   您将在此处完成中英文摘要的填写                   %?
%?                                                               %?
%? ************************************************************* %?

% ====================== 摘要 ======================

\chapter*{\texorpdfstring{\centering \heiti \sanhao \textbf{摘\quad 要}}{摘要}}
\addcontentsline{toc}{chapter}{\texorpdfstring{\sihao{摘\qquad 要}}{摘要}}

\linespread{1.5}\selectfont
\fangsong\xiaosi

\vspace{1em}

\xfabstract{

    为提升仓储搬运自动化水平，本文对智能仓储场景下的举升式AGV展开研究。通过三种驱动方案对比，选定双轮差速驱动方案。机械层面完成了底盘、升降机构和罩壳设计，并对支撑梁进行有限元验证。电气层面选定STM32F103C8T6为主控芯片，完成了电机选型和传感器接口电路设计。

    运动学层面推导了差速底盘方程，对四类典型轨迹进行仿真。路径规划层面采用Hybrid A*算法生成无碰撞路径，长度10.71\,m。循迹控制层面以5路红外传感器阵列进行PID仿真，RMS横向误差0.8\,cm。运动控制层面设计了分层PID控制器，在扰动条件下验证了跟踪与抗扰能力。

    仿真结果表明，该方案在结构、规划与控制层面具备可行性，可为样机研制提供参考。

}

\noindent{\heiti\xiaosi\textbf{关键词：}}
\fangsong{举升式AGV；智能仓储；差速驱动；路径规划；Hybrid A*；PID循迹控制}

% ====================== Abstract ======================

\chapter*{}
\vspace*{-1.15cm}
{\centering\bfseries\sanhao ABSTRACT\par}
\addcontentsline{toc}{chapter}{\sihao{ABSTRACT}}
\vspace{0.5cm}

\linespread{1.5}\selectfont
\rmfamily\xiaosi

\xfenabstract{
    With the rapid growth of e-commerce and intelligent logistics, traditional warehouse handling faces increasing challenges in efficiency and cost. To improve automation in warehouse material handling, this thesis investigates a lifting-type AGV for intelligent warehousing.

    First, by comparing three drive schemes, a dual-wheel differential drive architecture is adopted for the lifting AGV. For the mechanical structure, the chassis, drive wheels, lifting mechanism, and shell are designed, and finite element analysis is performed on the supporting beam. For the electrical system, the STM32F103C8T6 microcontroller is selected as the main controller, the servo motor is matched with a planetary reducer, and interface circuits are designed for infrared and ultrasonic sensors.

    In terms of kinematics, the differential-drive chassis kinematic equations are derived, and MATLAB simulations are conducted for four typical trajectories: straight, circular, S-shaped, and figure-eight. For path planning, the Hybrid A* algorithm is employed to generate collision-free paths in a warehouse grid map, yielding a path length of 10.71\,m. For line-following control, PID simulation with a 5-sensor infrared array achieves an RMS lateral error of 0.8\,cm. Finally, a dual-wheel PID controller is tested under disturbance conditions, demonstrating its disturbance rejection capability.

    Simulation results show that the proposed AGV design is feasible in structural design, path planning, and motion control, providing a reference for prototype development.
}

\vspace{1em}

\noindent\textbf{Key words:} lifting AGV; intelligent warehousing; differential drive; path planning; Hybrid A*; PID line-following control

\setcounter{secnumdepth}{2}
